\section{微分方程式}\label{節:微分方程式}
これまでは,関数が与えられたときに導関数を求めた.逆に,導関数についての条件から関数を決めることはできるだろうか.たとえば $k$ を定数として
\[
 y'=ky
\]
を満たす関数$y$を考えよう.このように,未知の関数とその導関数を含む方程式を\textbf{微分方程式}という.

$y=\e^{kx}$ はこの式を満たす.定数倍した$C\e^{kx}$も満たすが,ほかの形の解はないだろうか.

もし $y=C\e^{kx}$ なら,\,$y$ を既に見つけた解 $\e^{kx}$ で割った商 $\e^{-kx}y$ は定数 $C$ になるはずである.そこで,どの解 $y$ についても,この商が定数になることを示そう.積の微分と $y'=ky$ より
\begin{align*}
 \frac{\dd}{\dd x}(\e^{-kx}y)
 &=-k\e^{-kx}y+\e^{-kx}y'\\
 &=\e^{-kx}(y'-ky)\\
 &=0
\end{align*}
となる.定理\ref{定理:導関数同じ}より,導関数が$0$の関数は1つの区間で定数である.したがって$\e^{-kx}y=C$であり,
\[
 y=C\e^{kx}
\]
となる.ほかの形の解はないのだ.さらに$y(0)=A$も与えられれば,\,$C=A$ として解は1つに定まる.

\begin{bookpoint}
\textbf{POINT}\par
既知の解 $\e^{kx}$ で割った $\e^{-kx}y$ を微分して $0$ を示し,定数 $C$ と置いて $y=C\e^{kx}$ とする.
\end{bookpoint}

\paragraphbf{等比数列との関係}
解 $y=C\e^{kx}$ を,等しい間隔$h$ごとに取り出してみよう.\,$a_n=y(nh)$と置くと,
\begin{align*}
 a_{n+1}&=C\e^{k(n+1)h}\\
 &=\e^{kh}\cdot C\e^{knh}\\
 &=\e^{kh}a_n
\end{align*}
となる.つまり,\,$y'=ky$ の解を等間隔で見ると,公比 $\e^{kh}$ の等比数列になっている.等比数列の値を,その間の点まで指数関数でつないだものが $y'=ky$ の解である.

では逆に,解を知らないことにして,微分方程式 $y'=ky$ から等比数列を作れないだろうか.導関数は差商の極限であった.そこで,極限をとる前の差商で導関数を置き換えてみよう.間隔 $h$ を1つ決め,\,$y(nh)$ に当たる値を $a_n$ と書く.\,$y'=ky$ の左辺を差商 $\frac{a_{n+1}-a_n}{h}$ に,右辺の $y$ を $a_n$ に置き換えると,
\[
 \frac{a_{n+1}-a_n}{h}=ka_n
\]
となる.両辺に $h$ を掛けて $a_n$ を移項すれば,
\[
 a_{n+1}=(1+kh)a_n
\]
を得る.これは,項に一定の数 $p=1+kh$ を掛けて次の項を作る漸化式 $a_{n+1}=pa_n$ であり,\,$a_n$ は公比 $1+kh$ の等比数列になる.初めの値を $a_0=C$ とすれば,
\[
 a_n=C(1+kh)^n
\]
である.

この等比数列は,間隔 $h$ を細かくすると,解 $C\e^{kx}$ に近づくだろうか.実数 $x$ を1つ固定し,\,$0$ から $x$ までを $n$ 等分して $h=\frac xn$ と置く.\footnote{$x<0$ のときは $h<0$ であり,負の向きへ進むことになるが,差商の式は変わらない.}すると,\,$n$ 回進んだ $a_n$ が $y(x)$ に当たる値であり,
\[
 a_n=C\left(1+\frac{kx}{n}\right)^n
\]
となる.\,$n$ を大きくすることが,間隔を細かくすることに当たる.そこで,\,$n\to\infty$ のときの極限を求めよう.

$kx=0$ なら,すべての $n$ で $a_n=C$ であり,これは $C\e^0$ に等しい.以下,\,$kx\ne0$ とする.\,$n$ が $|kx|$ より大きければ $1+\frac{kx}{n}>0$ だから,そのような $n$ だけを考えれば,対数をとることができる.\,$t=\frac{kx}{n}$ と置くと,\,$n=\frac{kx}{t}$ より
\begin{align*}
 \log\left(1+\frac{kx}{n}\right)^n
 &=n\log\left(1+\frac{kx}{n}\right)\\
 &=kx\cdot\frac{\log(1+t)}{t}
\end{align*}
となる.\,$n\to\infty$ のとき $t\to0$ である.ここで $\log1=0$ に注意すると,
\[
 \frac{\log(1+t)}{t}=\frac{\log(1+t)-\log1}{t}
\]
は,\,$\log x$ の $x=1$ における差商にほかならない.\,$(\log x)'=\frac1x$ だから,その極限は $x=1$ での微分係数 $1$ である.したがって
\[
 \log\left(1+\frac{kx}{n}\right)^n\longrightarrow kx\qquad(n\to\infty)
\]
となる.指数関数は連続なので,両辺を $\e$ の指数に戻して
\[
 \lim_{n\to\infty}C\left(1+\frac{kx}{n}\right)^n=C\e^{kx}
\]
を得る.\,$C=1$,\,$kx=1$ の場合は,\,$\e$ を定めた極限 $\left(1+\frac1n\right)^n\to\e$ そのものである.

こうして,微分方程式 $y'=ky$ と等比数列の漸化式 $a_{n+1}=pa_n$ は,両方の向きに結び付いた.解を等間隔で見れば等比数列になり,逆に,細かく区切って作った等比数列の極限は指数関数になるのだ.指数関数は,いわば連続的な等比数列なのである.漸化式の側から見たこの対応は,『漸化式の超体系的解説』でも扱っている.

\paragraphbf{$y''=-y$ の場合}
もう1つ,\,$y''=-y$を考えてみよう.\,$\sin x,\cos x$はいずれも2回微分すると符号が変わるので,
\[
 y=A\cos x+B\sin x
\]
は解であり,\,$y(0)=A$,\,$y'(0)=B$を満たす.これ以外の解がないことを確かめよう.同じ初期条件をもつ別の解との差を$z$とすると,
\[
 z''=-z,\qquad z(0)=z'(0)=0
\]
である.ここでは少々発想しにくい式を使うので,結果からお見せしよう.\,$z^2+(z')^2$ を微分してやると
\begin{align*}
 \{z^2+(z')^2\}'
 &=2zz'+2z'z''\\
 &=2z'(z+z'')\\
 &=0
\end{align*}
となる.したがって$z^2+(z')^2$は定数で,\,$x=0$で$0$だから常に$0$である.2つの平方の和が$0$なので,\,$z=0$を得る.微分方程式の一般的な理論を使わずに,解の形と一意性を確かめられた.

この式は,\,$z''=-z$ の両辺に $z'$ を掛けると見つかる.
\[
 z'z''=-zz'
\]
の左辺は $\frac12(z')^2$ の導関数,右辺は $-\frac12z^2$ の導関数である.右辺を移項して2つを同じ側に集めると,\,$\frac12\{z^2+(z')^2\}$ の導関数が $0$ という式になるのだ.

時刻$t$での位置を$s(t)$とすると,速度は$s'(t)$であり,速度の変化率$s''(t)$を\textbf{加速度}という.たとえば$s''(t)=-s(t)$は,原点から離れるほど大きく,原点へ向かう加速度が生じる運動を表す.先ほどの結果から,位置は$s(t)=A\cos t+B\sin t$となり,周期的に往復する.このような運動を\textbf{単振動}という.同じ微分方程式は,変数を時刻と考えない場合にも,2回微分したときの関係によって三角関数を特徴付けている.

ここまでは,微分から分かることを調べた.次章では,導入で扱った面積を一般の関数について定め,その変化率が元の関数になることを示す.そこから,微分の逆によって積分を計算する方法を考えよう.
